\documentclass{article}
\usepackage[utf8]{inputenc}
\usepackage[T1]{fontenc}
\usepackage{amsmath}
\usepackage{amsfonts}
\usepackage{amssymb}
\usepackage[dvipsnames]{xcolor}
\usepackage{enumitem}
\usepackage{titlesec}
\usepackage{graphicx}
\usepackage[total={6.5in, 8.6in}, heightrounded]{geometry}
\usepackage{hyperref}

\hypersetup
{
	colorlinks = true,
	allcolors = OliveGreen
}

\graphicspath{{graphics/}}
\setenumerate[0]{label=\alph*)}
\setlength{\parindent}{0pt}
\setlength{\parskip}{8pt}
\setlength\fboxsep{0pt}
\renewcommand{\baselinestretch}{1.6}
\titleformat{\section}
{\normalfont \Large \bfseries \centering}{}{0pt}{}

\newcommand{\s}[1]{{\color{violet} #1}}



\begin{document}\Large Name: \rule{2in}{0.15mm} \hfill Problem Set 2 | Math 1180 | Cruz Godar \vspace{4pt}  \normalsize

\textit{Due Thursday, September 17th at 11:59 PM}

Complete the following problems and submit them as a PDF to Gradescope. You should show enough work that there is no question about the mathematical process used to obtain your answers, and so that your peers in the class could easily follow along. I encourage you to collaborate with your classmates, so long as you write up your solutions independently. Using AI tools to assist with the homework is not permitted. I write most of my problems to not require calculators, but you may find some of the problems require them --- that's perfectly fine! All exam problems will be written without the need for calculators.

~\\

1. The following graph plots these four equations in some order:

\begin{enumerate}

	\item $\displaystyle z = \frac{1}{5}\left( x^3 + y^3 \right)$.

	\item $\displaystyle z = xy$.

	\item $\displaystyle z = \sin(x) + \cos(y)$.

	\item $\displaystyle z = \sqrt{1 - x^2 - \left( \frac{y}{2} \right)^2}$.

\end{enumerate}

Toggle each on using the expressions in the sidebar and match each function to its graph (you can technically do this by just typing the functions in as new expressions, but the point is to analyze the graphs).

~\\

2. Find an equation for a cylinder in $\mathbb{R}^3$ parallel to the $y$-axis, with radius $4$, and whose central axis contains the point $(2, 3, 1)$.

~\\

3. Let $\vec{v} = \left< 0, 3, 4 \right>$ and let $\vec{w} = \vec{i} - 2\vec{j} + \vec{k}$.

\begin{enumerate}

	\item Find $\left| \left| \vec{v} \right| \right|$ and $\left| \left| \vec{w} \right| \right|$.

	\item Find $\vec{v} + \vec{w}$ and $2\vec{v} - \vec{w}$ in component form and sketch them both.

	\item Find unit vectors in the same direction as $\vec{v}$ and $\vec{w}$, respectively.

\end{enumerate}

~\\

In $\mathbb{R}^2$, two distinct lines are either parallel (if they have the same slope) or they intersect. In $\mathbb{R}^3$, there is a third possibility: nonparallel lines may not intersect, in which case they are called \textbf{skew}. In questions 4--6, determine whether the given pair of lines is parallel, intersecting, or skew.

4. $\left< -1 + 2t, 2 + 3t, 7t \right>$ and $\left< 3 - 6t, 1 - 9t, 10 - 21t \right>$.

5. $\left< 7 + 2t, -10 - t, 2 + t \right>$ and $\left< -8 + 3t, -3 - t, -9 + 5t \right>$.

6. $\left< 1 + t, 2 + t, 3 + t \right>$ and $\left< 4 + t, 5 - t, 7 + 2t \right>$.

~\\

7. Find an equation of a line passing through the points $(5, 1, 2)$ and $(3, 4, 1)$.

8. Find an equation of the plane containing the points $(1, 0, 1)$, $(-3, 1, 1)$, and $(2, 4, 0)$.

9. Find the intersection of the line from question 7 and the plane from question 8 if they exist, or show they don't intersect.

~\\

In questions 10--12, determine whether the statement is true or false, where all of the vectors are in $\mathbb{R}^3$. If the statement is true, briefly justify why. If it is false, provide a demonstrating example.

10. If $\vec{v} \bullet \vec{w} = 0$, then either $\vec{v} = \vec{0}$ or $\vec{w} = \vec{0}$.

11. If $\vec{v} \times \vec{w} = \vec{0}$, then either $\vec{v} = \vec{0}$ or $\vec{w} = \vec{0}$.

12. If $\vec{u} \bullet \left( \vec{v} \times \vec{w} \right) = 0$, then $\vec{u}$ is in the plane containing $\vec{v}$ and $\vec{w}$ when all three are drawn starting from the origin.

~\\

Let's try to find the general form of a unit vector in $\mathbb{R}^3$: rather than describing it as a vector $\vec{u} = \left< u_1, u_2, u_3 \right>$ with $u_1^2 + u_2^2 + u_3^2 = 1$, it can be useful to have two parameters that we can freely vary instead of three parameters with a restriction. We'll find such a form in the next two questions.

13. Suppose $u_3 = 0$, so that if $\vec{u}$ is drawn with its tail at the origin, then the tip of $\vec{u}$ is in the $xy$-plane, as in the graph below. If $\theta$ is its angle measured counterclockwise from the positive $x$-axis, drawn in red, find a general formula for $\vec{u}$ in terms of its angle $\theta$.

14. Now suppose $u_3 \neq 0$. We can think of $\vec{u}$ as being rotated up or down from the $xy$-plane by some angle $\varphi$; change the perspective on the graph and adjust $\varphi_0$ to see this in action. Hold $\theta$ fixed and look at the half-disk from the side, so that the problem effectively becomes 2D again. Use your answer to part a) and this perspective to find an expression for $\vec{u}$ in terms of $\theta$ and $\varphi$.


\end{document}